% Unidad R567: introducción, líneas francesas 37--68, pp. 6--7.
% Traducción: OpenAI Codex, GPT-6.1 Sol, esfuerzo Ultra.
A título informativo, damos a continuación el plan general previsto para este tratado, sujeto, por lo demás, a modificaciones posteriores, sobre todo en lo que respecta a los últimos capítulos:

\medskip
\begin{tabular}{rl}
  Capítulo primero. & --- El lenguaje de los esquemas.\\
  --- \textbf{II}. & --- Estudio global elemental de algunas clases de morfismos.\\
  --- \textbf{III}. & --- Cohomología de los haces algebraicos coherentes. Aplicaciones.\\
  --- \textbf{IV}. & --- Estudio local de los morfismos.\\
  --- \textbf{V}. & --- Procedimientos elementales de construcción de esquemas.\\
  --- \textbf{VI}. & --- Técnica de descenso. Método general de construcción de esquemas.\\
  --- \textbf{VII}. & --- Esquemas de grupos, espacios fibrados principales.\\
  --- \textbf{VIII}. & --- Estudio diferencial de los espacios fibrados.\\
  --- \textbf{IX}. & --- El grupo fundamental.\\
  --- \textbf{X}. & --- Residuos y dualidad.\\
  --- \textbf{XI}. & --- Teorías de intersección, clases de Chern, teorema de Riemann-Roch.\\
  --- \textbf{XII}. & --- Esquemas abelianos y esquemas de Picard.\\
  --- \textbf{XIII}. & --- Cohomología de Weil.\\
\end{tabular}
\medskip

En principio, todos los capítulos se consideran abiertos a ampliaciones, y siempre podrán añadírseles párrafos suplementarios más adelante; estos párrafos aparecerán en fascículos separados, para reducir los inconvenientes del modo de publicación adoptado.
Cuando uno de estos párrafos esté previsto o en preparación al publicarse un capítulo, se mencionará en el sumario de ese capítulo, aunque determinadas prioridades de urgencia hagan que su publicación efectiva tenga lugar mucho más tarde.
Para comodidad del lector, ofrecemos en un «Capítulo~0» diversos complementos de álgebra conmutativa, álgebra homológica y teoría de haces, utilizados en los capítulos de este tratado, que son más o menos conocidos, pero para los cuales no ha sido posible dar referencias cómodas de consultar.
Se recomienda al lector acudir al capítulo~0 solo mientras lee el tratado propiamente dicho, y en la medida en que los resultados a los que remitimos no le sean

\oldpage[I]{7}
suficientemente familiares.
Pensamos, además, que así la lectura de este tratado podrá ser para el principiante un buen método de familiarizarse con el álgebra conmutativa y el álgebra homológica, cuyo estudio, cuando no se acompaña de aplicaciones tangibles, resulta tedioso, e incluso deprimente, para un número considerable de personas.

\bigskip
\begin{center}*\quad*\quad*\end{center}
\bigskip
